% ---------------------------------------------------------------------
%  NỘI DUNG CHƯƠNG 1 — CƠ SỞ LÝ THUYẾT VÀ CÔNG NGHỆ
% ---------------------------------------------------------------------
\section{Cơ sở lý thuyết và công nghệ}
\label{sec:coso}

\begin{muctieu}
Phần này giới thiệu quy trình sinh câu hỏi từ tài liệu, các thành phần xử
lý dữ liệu và công nghệ thực tế trong hệ thống.
\end{muctieu}

\subsection{Kiến thức nền tảng}

Trong phạm vi đồ án, AI được dùng để hỗ trợ suy luận và sinh nội dung từ
học liệu qua mô hình ngôn ngữ, không phải để huấn luyện một bộ phân lớp mới.
Các khái niệm về tác tử và hệ thống trí tuệ nhân tạo được tham chiếu từ giáo
trình nền tảng~\cite{RussellNorvig2021}.

\subsubsection{Sinh câu hỏi có điều kiện theo nguồn}
Sinh câu hỏi từ tài liệu là quá trình chuyển nội dung văn bản thành câu hỏi
và đáp án dựa trên các ý trong nguồn. Đầu vào của AI ExamHub DAU là nội dung
đã trích xuất cùng môn học, số lượng, loại câu và phân bố độ khó/mức độ tư
duy. Kết quả được chuẩn hoá thành các trường như nội dung câu hỏi, lựa chọn,
đáp án, giải thích và đoạn bằng chứng. Gắn bằng chứng nguồn giúp giảng viên
đối chiếu đề xuất với học liệu trước khi lưu vào ngân hàng.

Khi có cấu hình khóa OpenAI, backend gửi yêu cầu đến dịch vụ mô hình ngôn
ngữ. Đây là suy luận qua API, không phải mô hình do đồ án huấn luyện. Khi
không có khóa hoặc phản hồi không dùng được, backend chuyển sang bộ sinh dự
phòng dựa trên các câu/đoạn trong tài liệu. Bộ dự phòng giúp luồng có thể
hoạt động mà không phụ thuộc dịch vụ ngoài, nhưng không đảm bảo chất lượng
ngữ nghĩa tương đương người biên soạn.

\subsubsection{Mô hình ngôn ngữ và sinh theo chỉ dẫn}
Mô hình ngôn ngữ ước lượng chuỗi từ tiếp theo dựa trên ngữ cảnh đã nhận.
Nhờ huấn luyện trên lượng văn bản lớn, một mô hình có thể làm theo chỉ dẫn
và tạo đầu ra với nhiều dạng cấu trúc; nghiên cứu về mô hình ngôn ngữ few-shot
mô tả khả năng điều kiện hoá đầu ra từ ví dụ trong ngữ cảnh
\cite{Brown2020}. Khả năng tạo câu tự nhiên không tự chứng minh nội dung là
đúng. Mô hình có thể diễn đạt trôi chảy nhưng suy luận sai, bỏ qua điều kiện
hoặc tự bổ sung chi tiết không có căn cứ.

Vì vậy, yêu cầu gửi đến mô hình cần làm rõ vai trò, nội dung nguồn, ma trận
câu hỏi, định dạng đầu ra, quy tắc đáp án và yêu cầu trích nguyên văn đoạn
bằng chứng. Ràng buộc cấu trúc giúp backend có thể phân tích JSON và từ
chối đầu ra sai định dạng, nhưng một JSON hợp lệ vẫn có thể chứa kiến thức
sai. Quy trình sản xuất cần có bước kiểm duyệt chuyên môn trước khi phát
hành.

\subsubsection{Sinh có truy xuất và khả năng truy nguyên}
Một hướng nghiên cứu liên quan là kết hợp mô hình sinh với truy xuất tri
thức bên ngoài. Retrieval-Augmented Generation (RAG) kết nối bộ sinh với
kho tri thức để cung cấp thông tin và dấu vết nguồn cho câu trả lời
\cite{Lewis2020RAG}. Trong nguyên mẫu hiện tại, backend đưa nội dung tài
liệu được người dùng gửi trực tiếp vào yêu cầu, rồi yêu cầu trả
\code{sourceEvidence}; dự án chưa có kho embedding/vector, bộ truy hồi theo
ngữ nghĩa hay pipeline RAG độc lập. Do đó, thuật ngữ “có căn cứ theo tài
liệu” ở đây mô tả điều kiện prompt và trích dẫn nguồn, không khẳng định
kiến trúc RAG đã được triển khai.

\subsubsection{Độ khó và mức độ nhận thức}
Độ khó của câu hỏi không chỉ phụ thuộc độ dài câu. Câu hỏi nhận biết/hiểu
thường yêu cầu giải thích hoặc nhận diện khái niệm; câu vận dụng yêu cầu
dùng quy tắc trong tình huống. Hệ thống lưu mức độ nhận thức
\code{understanding} hoặc \code{application} và các nhãn độ khó, đồng thời
cho phép giảng viên đặt số lượng mong muốn. Các nhãn trong metadata là đầu
vào của quá trình chọn câu hỏi, không phải kết quả đã được hiệu chuẩn theo
một thang đo tâm lý lượng.

\subsubsection{Chất lượng câu hỏi trắc nghiệm}
Câu trắc nghiệm cần có một đáp án đúng duy nhất, phương án nhiễu cùng phạm
trù và không vô lý, lời dẫn rõ ràng, tránh manh mối hình thức và tránh phụ
thuộc vào kiến thức ngoài phạm vi môn học. Bộ sinh của dự án kiểm tra một
số ràng buộc như số lượng/tính khác nhau của lựa chọn và khả năng đối chiếu
bằng chứng. Những kiểm tra này không xác định đầy đủ tính đúng đắn chuyên
môn hay chất lượng sư phạm, nên việc duyệt của giảng viên vẫn là một thành
phần bắt buộc của quy trình sử dụng an toàn.

\subsubsection{Trích xuất nội dung và OCR}
PDF được đọc bằng \code{pdf-parse}; DOC/DOCX được chuyển qua Mammoth, vốn
được thiết kế chủ yếu cho DOCX; TXT được đọc theo UTF-8. Backend cũng có
nhánh OCR ảnh dùng Tesseract.js, khởi tạo với dữ liệu ngôn ngữ tiếng Anh.
Giao diện hiện chỉ cho tải PDF, DOC/DOCX và TXT nên nhánh ảnh chưa phải
luồng tải lên thông thường. Văn bản quét, DOC không tương thích, bố cục phức
tạp hoặc tài liệu tiếng Việt đi qua bộ nhận dạng tiếng Anh có thể làm sai
ký tự và ranh giới câu; nội dung trích xuất cần được rà soát trước khi dùng.

\subsubsection{Mammoth và Tesseract.js}
Mammoth tập trung vào việc chuyển DOCX sang HTML/văn bản, còn Tesseract.js
cung cấp giao diện OCR JavaScript cho ảnh. Việc có thư viện trong backend
không đồng nghĩa chất lượng trích xuất tiếng Việt đã được bảo đảm; chất
lượng cần được đo riêng trên tập tài liệu phù hợp
\cite{MammothDocs,TesseractDocs}.

\subsubsection{Kiểm tra kết quả sinh}
Kết quả mô hình không được đưa thẳng vào đề. Backend chuẩn hoá và kiểm tra
các trường, loại câu, số lựa chọn, tính duy nhất của lựa chọn và đáp án,
mức độ tư duy, độ khó, đồng thời đối chiếu bằng chứng với văn bản nguồn. Nếu
kết quả không đạt điều kiện, hệ thống dùng kết quả dự phòng hoặc loại câu
không hợp lệ. Đây là kiểm tra cấu trúc và một số ràng buộc nội dung; nó
không thay thế việc giảng viên thẩm định kiến thức.

\subsubsection{Mô hình rủi ro và trách nhiệm sử dụng}
Khung quản trị rủi ro AI của NIST tổ chức hoạt động quanh các chức năng
quản trị, lập bản đồ, đo lường và quản lý rủi ro
\cite{NISTAIRMF}. Hướng dẫn UNESCO về AI tạo sinh trong giáo dục nhấn mạnh
cách tiếp cận lấy con người làm trung tâm, độ tuổi phù hợp và bảo vệ dữ
liệu cá nhân \cite{UNESCOGenAI}. Từ góc nhìn hệ thống, rủi ro chính gồm câu
hỏi sai đáp án, dữ liệu nguồn nhạy cảm bị gửi ra ngoài, OCR sai, lạm dụng
tài khoản và phân quyền sai. Biện pháp đề xuất là giảm thiểu dữ liệu gửi
đi, giới hạn quyền, không dùng đầu ra AI làm kết quả cuối cùng nếu chưa
được rà soát, và ghi rõ dịch vụ bên ngoài có thể nhận nội dung tài liệu.

\begin{figure}[H]
\centering
\resizebox{0.96\textwidth}{!}{%
\begin{tikzpicture}[node distance=8mm and 10mm]
  \node[objbox] (source) {Tài liệu\\ PDF/DOCX/TXT};
  \node[clsbox, right=of source] (extract) {Trích xuất\\ văn bản};
  \node[clsbox, right=of extract] (prompt) {Đặt yêu cầu\\ + ma trận đề};
  \node[clsbox, right=of prompt] (generate) {LLM tùy chọn\\ / sinh dự phòng};
  \node[clsbox, below=of generate] (validate) {Kiểm tra cấu trúc,\\ đáp án, bằng chứng};
  \node[objbox, left=of validate] (review) {Giảng viên\\ rà soát / sửa};
  \node[objbox, left=of review] (bank) {Ngân hàng\\ câu hỏi};
  \draw[flowarr] (source) -- (extract);
  \draw[flowarr] (extract) -- (prompt);
  \draw[flowarr] (prompt) -- (generate);
  \draw[flowarr] (generate) -- (validate);
  \draw[flowarr] (validate) -- (review);
  \draw[flowarr] (review) -- (bank);
\end{tikzpicture}}
\caption{Quy trình khai thác tài liệu và rà soát câu hỏi trong hệ thống.}
\label{fig:ai-pipeline}
\end{figure}
Hình~\ref{fig:ai-pipeline} mô tả luồng logic; kết quả của bước kiểm tra
tự động chỉ là cổng cấu trúc ban đầu. Phê duyệt nội dung vẫn thuộc trách
nhiệm giảng viên.

\subsection{Công cụ và thư viện}

\begin{table}[H]
\centering
\caption{Công nghệ được khai báo trong các manifest và cấu hình triển khai.}
\label{tab:congcu}
\begin{tabular*}{\textwidth}{@{\extracolsep{\fill}}>{\bfseries\color{daunavy}}p{0.20\textwidth}>{\color{daumaroon}}p{0.23\textwidth}p{0.47\textwidth}}
\toprule
\rowcolor{gray100}\bfseries\color{daunavy}Thành phần & \bfseries\color{daumaroon}Phiên bản & \bfseries\color{daunavy}Vai trò \\
\midrule
React & 18.3.1 & Xây dựng giao diện web \\
TypeScript & 5.6.3 & Kiểm tra kiểu cho mã giao diện \\
Vite & 5.4.11 & Phát triển và đóng gói frontend \\
Node.js & 20 (Docker) & Môi trường chạy backend \\
Express & 5.1.0 & Xây dựng REST API \\
PostgreSQL & 16 (Docker) & Lưu trữ dữ liệu quan hệ \\
Tesseract.js & 7.x & Thành phần OCR được khai báo trong backend \\
pdf-parse, Mammoth & Theo package.json & Đọc PDF và DOCX \\
Docker Compose & Theo môi trường cài đặt & Khởi chạy các container phát triển \\
\bottomrule
\end{tabular*}
\end{table}

\subsection{Tài liệu nền tảng và công nghệ liên quan}

Các mô hình ngôn ngữ hiện đại có thể sinh văn bản theo chỉ dẫn, nhưng đầu ra
có thể chứa thông tin sai hoặc không bám sát nguồn. Vì vậy, hệ thống ứng
dụng nguyên tắc có điều kiện theo tài liệu, yêu cầu bằng chứng và kiểm tra
đầu ra trước khi giảng viên duyệt; tài liệu API mô hình là nguồn tham khảo
cho cách tích hợp dịch vụ suy luận~\cite{OpenAIAPI}. Nền tảng còn sử dụng
các thành phần ứng dụng web và cơ sở dữ liệu đã được chuẩn hoá trong tài
liệu chính thức của React, Express và PostgreSQL~\cite{ReactDocs,TypeScriptDocs,ViteDocs,ExpressDocs,PostgreSQLDocs}.
Khác với một nghiên cứu mô hình học máy, đồ án tập trung vào tích hợp các
luồng nghiệp vụ và không đưa ra so sánh benchmark giữa các mô hình.

\subsubsection{Cơ sở dữ liệu quan hệ và tính toàn vẹn}
PostgreSQL hỗ trợ khoá chính, khoá ngoại, ràng buộc và giao dịch, phù hợp
với các quan hệ giữa tài khoản, đề, câu hỏi, lượt thi và bài nộp
\cite{PostgreSQLDocs}. Câu hỏi và câu trả lời được biểu diễn dưới dạng JSONB
ở một số cột để giữ lựa chọn và đáp án theo cấu trúc linh hoạt; những liên
kết cần truy vấn và toàn vẹn tham chiếu vẫn có bảng liên kết riêng. Cách
phân chia này tạo thuận lợi cho metadata, nhưng cần tránh để hai biểu diễn
\code{question\_ids} và \code{exam\_questions} sai khác nhau.

\subsubsection{Giao diện thành phần và hợp đồng HTTP}
React tổ chức giao diện thành các component có trạng thái cục bộ và props
\cite{ReactDocs}; TypeScript bổ sung kiểu dữ liệu tĩnh cho mô hình người
dùng, câu hỏi, đề và lượt thi \cite{TypeScriptDocs}. Vite cung cấp máy chủ
phát triển và quy trình build tài nguyên frontend
\cite{ViteDocs}. Backend dùng Node.js và Express để khai báo middleware,
xác thực và tuyến HTTP \cite{NodeDocs,ExpressDocs}. Những công nghệ này
giảm công sức tích hợp, nhưng vẫn cần hợp đồng API nhất quán, xử lý lỗi và
quyền truy cập ở backend.

\subsubsection{Kiểm thử phần mềm và khả năng tái lập}
Bộ kiểm thử backend sử dụng module \code{node:test} có sẵn trong Node.js;
test runner cung cấp cách tổ chức và thực thi các test đồng bộ hoặc bất
đồng bộ \cite{NodeTestDocs}. Compose mô tả dịch vụ backend, frontend và cơ
sở dữ liệu trong một cấu hình để có thể dựng lại stack
\cite{DockerComposeDocs}. Kết quả kiểm thử chỉ có ý nghĩa đối với phiên bản
mã, cấu hình và môi trường đã chạy; vì vậy khi tái lập cần ghi lại commit,
phiên bản Node.js, biến môi trường đã che bí mật và tình trạng cơ sở dữ liệu.
